\documentclass[11pt]{article}
\usepackage{mathblatt}
% gesetzt von werkzeuge/setzer.py; Lösungen nach bau/bauregeln.md „Lösungen“.
\newcommand{\kennung}{QG4}
\begin{document}
\pfheftstil
\fancyfoot[C]{{\footnotesize\color{mbgrau}\kennung\hfill\thepage}}
\pfheftkopf{Bruch als Anteil}{Klasse 7 \textperiodcentered{} Lösungen}

\begin{pfloesung}{}
\lz{1b)}{$\frac{1}{6}$}{$\frac{1}{6}$: 1 von 6 gleich großen Teilen ist grau.}{}
\end{pfloesung}
\begin{pfloesung}{}
\lz{2.}{$\frac{7}{12}$}{$\frac{7}{12}$: 7 von 12 Kästchen sind grau (nicht 7 von 5: die weißen sind nur der Rest).}{}
\end{pfloesung}
\begin{pfloesung}{}
\lz{3.}{grau $\frac{4}{10}$, weiß $\frac{6}{10}$}{$\frac{4}{10}$; $\frac{6}{10}$: grau 4, weiß 6, zusammen 10 Teile.}{}
\end{pfloesung}
\begin{pfloesung}{}
\lz{4.}{$\frac{13}{24}$}{$\frac{13}{24}$: Mädchen $24 - 11 = 13$; das Ganze sind alle 24 Kinder.}{}
\end{pfloesung}
\begin{pfloesung}{}
\lz{5.}{In beiden Klassen gleich: je die Hälfte.}{Gleich groß. 7a: $\frac{12}{24} = \frac{1}{2}$; 7b: $\frac{10}{20} = \frac{1}{2}$. In beiden Klassen trägt genau die Hälfte ein rotes Shirt – mehr Kinder heißt nicht größerer Anteil.}{}
\end{pfloesung}
\begin{pfloesung}{}
\lz{6b)}{5 Teile}{5 Teile färben: Der Streifen hat 8 Teile, also ist jeder Teil ein Achtel.}{}
\end{pfloesung}
\begin{pfloesung}{}
\lz{7.}{4 Teile}{4 Teile färben: $6 : 3 = 2$ Teile sind ein Drittel, $2 \cdot 2 = 4$ Teile sind zwei Drittel.}{}
\end{pfloesung}
\begin{pfloesung}{}
\lz{8.}{20 Kästchen}{20 Kästchen schraffieren: $24 : 6 = 4$, $4 \cdot 5 = 20$.}{}
\end{pfloesung}
\begin{pfloesung}{}
\lz{9.}{1 von 4 gleichen Teilen}{Ein Viertel: das Quadrat in 4 gleich große Teile teilen (zwei Mittellinien oder beide Diagonalen) und einen Teil färben.}{}
\end{pfloesung}
\begin{pfloesung}{}
\lz{10.}{Ja, es reicht; 3 Felder bleiben frei.}{Ja; Tomaten $20 : 5 \cdot 3 = 12$ Felder, Salat $20 : 4 = 5$ Felder, zusammen $12 + 5 = 17$ Felder; frei bleiben $20 - 17 = 3$ Felder.}{}
\end{pfloesung}
\begin{pfloesung}{}
\lz{11b)}{$\frac{5}{16}$}{$\frac{5}{16}$: in halben Kästchen zählen – das Rechteck hat 16 halbe Kästchen, grau sind $2 \cdot 2 + 1 = 5$ halbe – oder $2\frac{1}{2}$ von 8 Kästchen, gleich viel.}{}
\end{pfloesung}
\begin{pfloesung}{}
\lz{12.}{$\frac{1}{6}$}{$\frac{1}{6}$: die großen Teile sind doppelt so groß wie die kleinen; in kleine Stücke geteilt hat der Kreis $2 + 2 + 1 + 1 = 6$ Stücke, grau ist 1.}{}
\end{pfloesung}
\begin{pfloesung}{}
\lz{13.}{$\frac{3}{8}$}{kleinster Sektor ein Achtel; der Kreis hat 8 Achtel, markiert sind $2 + 1 = 3$ Achtel.}{}
\end{pfloesung}
\begin{pfloesung}{}
\lz{14.}{$\frac{2}{8}$ $= \frac{1}{4}$}{$\frac{1}{4}$: grau sind 1 ganzes Kästchen und 2 halbe, zusammen 2 Kästchen von 8; $\frac{2}{8} = \frac{1}{4}$.}{}
\end{pfloesung}
\begin{pfloesung}{}
\lz{15.}{$\frac{3}{12}$ $= \frac{1}{4}$, genau ein Viertel}{$\frac{3}{12}$, genau ein Viertel: in halben Stücken gezählt hat die Pizza 12 Stücke, Ben isst $2 + 1 = 3$; $\frac{3}{12} = \frac{1}{4}$.}{}
\end{pfloesung}
\begin{pfloesung}{}
\lz{16b)}{15\,min}{15\,min: $40 : 8 = 5$; $5 \cdot 3 = 15$.}{}
\end{pfloesung}
\begin{pfloesung}{}
\lz{17.}{$1{,}5$\,l}{$1{,}5$\,l: $1{,}8 : 6 = 0{,}3$; $0{,}3 \cdot 5 = 1{,}5$.}{}
\end{pfloesung}
\begin{pfloesung}{}
\lz{18.}{32\,l}{32\,l: gefüllt $80 : 5 \cdot 3 = 48$\,l; es fehlen $80 - 48 = 32$\,l (gefragt ist der Rest, nicht die Füllmenge).}{}
\end{pfloesung}
\begin{pfloesung}{}
\lz{19.}{1\,750\,m}{1\,750\,m: $2{,}5$\,km $= 2\,500$\,m; $2\,500 : 10 = 250$; $250 \cdot 7 = 1\,750$.}{}
\end{pfloesung}
\begin{pfloesung}{}
\lz{20.}{Ja: Rest $0{,}25$\,l, genau ein Viertelliter.}{Ja, genau: $0{,}75 : 3 = 0{,}25$; $0{,}25 \cdot 2 = 0{,}5$\,l getrunken; Rest $0{,}75 - 0{,}5 = 0{,}25$\,l, das ist genau ein Viertelliter.}{}
\end{pfloesung}
\begin{pfloesung}{}
\lz{21b)}{$\frac{5}{10}$ $= \frac{1}{2} = 50\,\%$}{$\frac{5}{10} = \frac{1}{2} = 50\,\%$: 5 von 10 Teilen, das ist die Hälfte.}{}
\end{pfloesung}
\begin{pfloesung}{}
\lz{22.}{6 Kästchen}{6 Kästchen markieren: $20\,\% = \frac{1}{5}$; $30 : 5 = 6$.}{}
\end{pfloesung}
\begin{pfloesung}{}
\lz{23.}{$\frac{15}{20}$ $= \frac{3}{4} = 75\,\%$}{$\frac{15}{20} = \frac{3}{4} = 75\,\%$: $20 : 4 = 5$ Kästchen sind ein Viertel, 15 Kästchen sind drei Viertel.}{}
\end{pfloesung}
\begin{pfloesung}{}
\lz{24.}{2 Kästchen}{2 Kästchen markieren: $10\,\% = \frac{1}{10}$; $20 : 10 = 2$.}{}
\end{pfloesung}
\begin{pfloesung}{}
\lz{25.}{Nein: $75\,\%$ belegt, unter $80\,\%$.}{Nein: $\frac{30}{40} = \frac{3}{4} = 75\,\%$ ($40 : 4 = 10$ Plätze sind ein Viertel, 30 Plätze drei Viertel); $75\,\% < 80\,\%$.}{}
\end{pfloesung}
\begin{pfloesung}{}
\lz{26.}{$\frac{8}{20}$ $= \frac{2}{5} = 40\,\%$}{$\frac{8}{20} = \frac{2}{5} = 40\,\%$: $20 : 5 = 4$ Kästchen sind ein Fünftel ($20\,\%$), 8 Kästchen sind zwei Fünftel.}{}
\end{pfloesung}
\begin{pfloesung}{}
\lz{27.}{$\frac{4}{8}$}{kleinster Sektor ein Achtel; markiert sind $2 + 1 + 1 = 4$ Achtel – die Hälfte.}{}
\end{pfloesung}
\begin{pfloesung}{}
\lz{28.}{15 Kästchen}{15 Kästchen schraffieren: $25 : 5 = 5$, $5 \cdot 3 = 15$.}{}
\end{pfloesung}
\begin{pfloesung}{}
\lz{29.}{150\,l}{150\,l: gefüllt $900 : 6 \cdot 5 = 750$\,l; es passen noch $900 - 750 = 150$\,l hinein.}{}
\end{pfloesung}
\begin{pfloesung}{}
\lz{30.}{2\,000\,g}{2\,000\,g: $3{,}2$\,kg $= 3\,200$\,g; $3\,200 : 8 = 400$; $400 \cdot 5 = 2\,000$.}{}
\end{pfloesung}
\begin{pfloesung}{}
\lz{31.}{Nein: 290\,l, es bleiben 10\,l Platz.}{Nein: gefüllt $300 : 5 \cdot 4 = 240$\,l; $240 + 50 = 290$\,l, das sind 10\,l weniger als 300\,l.}{}
\end{pfloesung}
\end{document}
